\section{Conclusion and Future Scope}\label{sec:conclusion}
We built SurroSelect, a procedure and open tool that chooses a stand-in for a slow simulation by cross-validation, reports ties, checks whether the POD basis limits accuracy and forecasts the error for more runs, and we tested it against criteria committed in advance on eight problems solved by an independent finite-element library. Selection was within a factor \SoneMax{} of the oracle (the best pool member) in all 24 cells, although one fixed full-field Gaussian process does almost as well; the basis check raised a flag only in cases where a full-field model was better, but missed most such cases; the tool did not beat the benchmark's selection among networks and Gaussian-process variants, and its forecast was usable from small budgets only. A 44-problem study showed that the winner depends on smoothness, the number of inputs, the budget, the training effort of the networks and how widely a kernel method's regularisation is searched, and that the simple rule ``reduced basis wins at large budgets'' fails on \TpNotClear{} of \TpN{} problems solved by third-party code.
Concrete next steps are (i) a pool that includes a tuned network, with a rule for when to spend its cost; (ii) a forecast that models the error floor and gives a calibrated interval; (iii) tests on real field data from simulations that are much slower than ours, with the full cost of the simulations included; (iv) a public benchmark such as PDEBench~\cite{takamoto2022}; and (v) independent human replication.
